\section{Conclusion}
In this paper, we investigated the optimal trade-off between parameters and compute per example for maximizing model capacity.
% Our experiments and analysis, underscores the role of sparsity, as a factor that controls the balance between total number of parameters and \flops per example, in the scaling laws for \moes, showing that the most efficient model configuration depends on balancing model size, training compute, and sparsity level. The optimal recipe for balancing \flops per example and parameter count in \moes depends on the objective as well other resource constraints. 
Our findings indicate that sparsity, as a knob that controls \flops per example in \moes, is a powerful mechanism for optimizing model performance under constrained training compute budgets. By balancing the total number of parameters, compute, and sparsity, \moes can be scaled more effectively. These insights provide valuable guidance for scaling language models, especially for \moes, where the trade-offs between parameters and \flops must be carefully managed.

\moes were originally introduced to allow increasing model capacity without a significant increase in inference cost. Our experiments show that under fixed total training compute budget increasing sparsity in \moes leads to smaller \flops per example, higher number of parameters, and lower pretraining loss simultaneously. In other words, in the context of \moes, if there are no constraints on the total number of parameters, increasing the capacity of the model through parameter count seem to be the optimal strategy if lower pretraining loss is the main goal. On the other hand, when comparing how well the pretraining performance transfers to various downstream tasks, denser models  exhibit better transfer performance on certain types of task that potentially rely on deeper processing of the input vs the knowledge stored in the parameters of the model. 
This potentially signals the importance of the role of \flops per example in increasing the capacity of the model during inference. 
Our experiments demonstrate that \moes use Chain-of-Thought prompting more effectively than dense models, achieving better performance when allocated additional computational resources during inference. 
This observation reveals an interesting direction to improve the performance efficiency of \moes at inference time.
% \todo[inline]{talk about CoT as an example approach that could make sparser \moes more efficient}
% Concretely, this may be studied by verifying whether\moes use Chain-of-Thought prompting more effectively than dense models at inference.


Future work will focus on determining the optimal balance between FLOPs per example and parameter count, with an emphasis on conducting in-depth analyses of model performance across diverse downstream tasks. A key direction will involve exploring strategies to balance parameter allocation and computational demands to minimize inference costs. Developing scaling law studies to identify optimal approaches for achieving efficiency and performance during inference represents a critical area for further investigation.

Another important avenue will be to examine how the findings on the role of sparsity in \moes generalize to architectures or approaches that employ different mechanisms for independently adjusting FLOPs per example and the number of trainable parameters. Additionally, an intriguing direction for future exploration is the study of scaling behaviors in models that enable negative sparsity values through parameter sharing.
